\section{A single stable spin factor with an arbitrary entire partner}
\label{sec:stable-spin-arbitrary-partner}

\paragraph{Status and attribution.}
HD-FT27--28, reserved question Q2ek, is a separate strengthening of
the single-factor case of HD-FT25--26.  Root proposed the uniform
trace, projected-entry count and cube-quotient argument;
\texttt{/root/g34\_eightfolds} reconstructed the full proof below.
The accepted inputs are HD-FT23--24, including every faithful spin
irreducible and compatible outer operator, and the effective actual
quotient-genus theorem HD-GC06.  No new external character or genus
estimate is assumed. Complete author, root and independent whole-proof
readings pass; their exact source, dependency and arithmetic scopes
are recorded separately.

For a finite-order operator $T$ on an $e$-dimensional complex space,
put $\nu(T)=e-\max_\zeta\dim\ker(T-\zeta I)$.
The actual moving rank $\operatorname{rank}(1-T)$ is at least
$\nu(T)$, but may be larger, in particular for a nonidentity scalar.

\subsection*{HD-FT27. A cubic trace bound with only finitely many projected entries}

\begin{proposition}[HD-FT27]
Let $C\longrightarrow\mathbf P^1$ be an actual connected finite
$J$-Galois cover, and let $V$ be a faithful complex coefficient
representation of $J$, of degree $e$.  Suppose there is an actual
surjection
\[
 x:J\longrightarrow P,\qquad P\in\{A_n,S_n\},\qquad
 n\ge10^6,\qquad e\ge n^3.
 \tag{SSP.1}
\]
Write $f(y)$ for the number of fixed natural letters of $y\in P$.
Assume that for every $g\in J$ the following two bounds hold:
\[
 \frac{|\operatorname{tr}_V(g)|}{e}
 \le\left(\frac{\max\{f(x(g)),1\}}n\right)^3,
 \qquad x(g)\ne1\ \Longrightarrow\quad
 \frac{|\operatorname{tr}_V(g)|}{e}\le\frac78.
 \tag{SSP.2}
\]
Then the genuine coefficient inertia tuple of this cover cannot
have moving-rank sum $D=2e+6$.  No lower bound on its number of
projected nonidentity entries is required.
\end{proposition}

\begin{proof}
Let $g_1,\ldots,g_r$ be the actual inertia tuple, put $T_i=V(g_i)$,
and let $k$ be the number of indices for which $x(g_i)\ne1$.
For any finite-order operator, its eigenvalues have absolute value
one.  If its largest eigenspace has dimension $m$, the triangle
inequality gives $|\operatorname{tr}T|\ge2m-e$.  Consequently
\[
 \nu(T)\ge\frac{e-|\operatorname{tr}T|}{2},\qquad
 x(g_i)\ne1\ \Longrightarrow\quad
 \operatorname{rank}(1-T_i)\ge\frac e{16}.
 \tag{SSP.3}
\]
If $k\ge33$, this would give
$D/e\ge33/16=2+1/16$, whereas $6/e<1/16$ because
$e\ge n^3>96$.  Therefore
\[
                              k\le32.
 \tag{SSP.4}
\]
Genuine entries killed by $x$ have not been discarded from $D$;
only their nonnegative contributions will be omitted in a lower
bound.

Let $\Omega=[n]^3$, with the coordinatewise action of $P$, and
write $I_\Omega(y)$ for the permutation index of $y$ on this
possibly intransitive set.  Fix one original operator $T_i$ of
order $L_i$.  The order of $x(g_i)$ divides $L_i$, since $V$ is
faithful and $x$ is an actual quotient of the coefficient image.
For every eigenvalue target $\zeta$ of $T_i$, the cyclic Fourier
formula and (SSP.2), applied to every original power, imply
\[
 \begin{split}
 \frac{\dim\ker(T_i-\zeta I)}e
 &\le\frac1{L_i}\sum_{j=0}^{L_i-1}
                \frac{|\operatorname{tr}(T_i^j)|}{e}\\
 &\le\frac1{L_i}\sum_{j=0}^{L_i-1}
       \left\{\frac{f(x(g_i)^j)^3}{n^3}
                    +\frac{\mathbf1_{f(x(g_i)^j)=0}}{n^3}\right\}.
 \end{split}
 \tag{SSP.5}
\]
The first average in the last line is the number of cycles of
$x(g_i)$ on $\Omega$, divided by $n^3$.  Indeed its fixed-point
count on $\Omega$ is exactly $f(x(g_i)^j)^3$, and averaging over
$L_i$ repeats each power of its natural cyclic quotient equally
often.  Thus, for $x(g_i)\ne1$,
\[
 \frac{\operatorname{rank}(1-T_i)}e
 \ge\frac{\nu(T_i)}e
 \ge\frac{I_\Omega(x(g_i))}{n^3}-\frac1{n^3}.
 \tag{SSP.6}
\]
For $x(g_i)=1$, its cube index is zero and the inequality with
zero error follows just from nonnegative actual moving rank.
Summing only the necessary errors gives
\[
 \frac De\ge\frac1{n^3}\sum_i I_\Omega(x(g_i))
                         -\frac{k}{n^3}.
 \tag{SSP.7}
\]
This uses the full original orders, including powers lying in
$\ker x$, and makes no assumption that such coefficient powers
are scalar.

The quotient $C/\ker x\longrightarrow\mathbf P^1$ is an actual
connected $P$-Galois cover.  The $P$-set $\Omega$ has exactly
five transitive orbits: the all-equal triples, the three specified
patterns with exactly two equal coordinates, and the triples with
three distinct coordinates.  Both $A_n$ and $S_n$ are transitive
on each of these sets in the stated range.  Let $Y_1,Y_2,Y_3$
be the actual connected natural, ordered-distinct-pair and
ordered-distinct-triple quotient curves.  Let $X_3$ be the
unordered-triple quotient.  There is an actual degree-six map
$Y_3\longrightarrow X_3$; Riemann--Hurwitz and HD-GC06 give
\[
 g(Y_3)\ge6g(X_3)-5>\frac{n^2}{48}-5.
 \tag{SSP.8}
\]
There are three copies of $Y_2$ among the five cube orbits.
Summing Riemann--Hurwitz over their actual quotient covers yields
\[
 \begin{split}
 \sum_i I_\Omega(x(g_i))
 &=2n^3-10+2\{g(Y_1)+3g(Y_2)+g(Y_3)\}\\
 &>2n^3+\frac{n^2}{24}-20.
 \end{split}
 \tag{SSP.9}
\]
The first two genera are nonnegative; no branch-list realization
or connectedness of the whole cube set is asserted.

Combining (SSP.4), (SSP.7) and (SSP.9) now gives
\[
 \frac De>2+\frac1{24n}-\frac{52}{n^3}
          >2+\frac1{25n}>2+\frac6e.
 \tag{SSP.10}
\]
The middle inequality is equivalent to $n^2>31200$.
The last follows from $e\ge n^3>150n$.  Both hold for
$n\ge10^6$.  This contradicts $D=2e+6$ and proves the proposition.
\end{proof}

\subsection*{HD-FT28. Every single stable spin factor, without maximality}

\begin{theorem}[HD-FT28]
Suppose an actual finite coefficient image $J$ preserves the
whole matrix algebra of one irreducible spin factor $W$ of
$\widehat A_n=2.A_n$, where $n\ge10^6$.  Its finite projective
image on $W$ contains the represented $A_n$ and normalizes it.
On the entire coefficient space write
\[
                    V=W\otimes R,\qquad
 a=\dim W,\quad b=\dim R\ge1,\quad e=ab.
 \tag{SSP.11}
\]
Here $R$ is the entire external tensor partner; no primitivity,
dimension bound, decomposition, character field or projective
support floor is assumed for that partner.  Then this stable
spin factor cannot occur in a connected sphere cover with actual
moving-rank sum $D=2e+6$.  In particular it cannot occur in an
actual rank-nine sixfold curve supplier.  Maximality of its
alternating natural degree is not required.
\end{theorem}

\begin{proof}
Stability of the factor algebra gives genuine finite projective
quotient actions on $W$ and $R$.  Conjugation on the represented
perfect spin group gives an actual homomorphism
$x:J\longrightarrow\operatorname{Aut}(A_n)=S_n$.
Its image contains $A_n$, so is $A_n$ or $S_n$.
There is no choice of natural permutation separately at each
branch point.  The complete perfect-core and compatible-odd
extension argument is the one in HD-FT24.

That argument also supplies one fixed finite double-cover model
for every compatible operator on $W$.  If no odd automorphism
stabilizes the spin irreducible, only the inner cover is needed;
otherwise it extends to the indicated $\widehat S_n$.
For any $g\in J$, take its representative $A_g$ in that finite
cover model, with natural projection $x(g)$, and absorb the
reciprocal scalar in the other tensor factor so that exactly
\[
                    V(g)=A_g\otimes B_g.
 \tag{SSP.12}
\]
This is a normalization of the two tensor lifts, not an assertion
that the two assignments are ordinary representations of $J$.
Let $L$ be the order of the original finite coefficient operator
$V(g)$.  Since $x(g)^L=1$, the cover element defining $A_g^L$
is central, so $A_g^L=\epsilon I$ with $\epsilon\in\{1,-1\}$.
Taking the $L$th power of (SSP.12) shows
$B_g^L=\epsilon^{-1}I$.  In particular $B_g$ is finite order
and every one of its eigenvalues has absolute value one.
For every integer $j$, therefore,
\[
 \frac{|\operatorname{tr}_V(g^j)|}{e}
 =\frac{|\operatorname{tr}(A_g^j)|}{a}
                 \frac{|\operatorname{tr}(B_g^j)|}{b}
 \le\frac{|\operatorname{tr}(A_g^j)|}{a}.
 \tag{SSP.13}
\]
Neither factor is required to have order $L$, and no power of
the entire partner is declared scalar except the displayed
$L$th-power consequence.

HD-FT23 gives the uniform $7/8$ bound whenever $x(g)^j\ne1$.
HD-FT24 gives $a\ge n^3$ and its cubic fixed-letter bound for
all these powers, including central ones, since
$n\ge10^6\ge65536$.  Thus (SSP.13) proves both hypotheses
(SSP.2) for the whole coefficient representation, and
$e\ge a\ge n^3$.  HD-FT27 gives the contradiction.

Finally, in the actual rank-nine inducing-source interface the
same degree $e>6$ forces the inducing base to be the sphere:
base genus at least two contradicts the multiplicity Euler
identity, and genus one has at most six genuine moving-rank
units and cannot have period rank nine.  On the sphere the
necessary budget is $D=2e+6$.  These are the already accepted
actual-cover conclusions of D-BI-18, also used explicitly in
HD-FT25.  The present sphere argument does not additionally
need the usual twelve-genuine-mark lower bound.
\end{proof}

\paragraph{Strict new scope and retained boundaries.}
HD-FT25 required the entire partner floor
$\nu_R(S)/b\ge1/n$ for every nonscalar partner operator;
HD-FT26 supplied it from maximality in a complete primitive
decomposition.  HD-FT28 removes that floor and maximality for
one stable spin factor, retaining an arbitrary entire partner
at the larger explicit threshold $10^6$.  No component selected
from a permuted spin layer is declared stable, and no new
statement about such permuted layers is proved here.  Exceptional
cover families at $A_6,A_7$ remain separate.  No original-field
descent, polarization comparison or source existence is inferred.
